% No preamble. Prefix fbt-. Conditional qualitative positive block transfer.
\section{Qualitative transfer of a positive leading block drift}
\label{fbt-section}

The compactness and moment-passage inputs in this section are
available for the actual homogeneous local macrostep. Its
positivity hypothesis is explicitly additional and unproved.
The conclusion requires no quantitative finite-volume rate.

\subsection{The actual fixed-length coefficients}
\label{fbt-setup}

Fix the primitive ratio, elementary composite ratio \(R\),
grouping \(J\), common analytic domains and extraction
convention. Fix a finite \(q\ge1\), chosen to be a multiple
of the activity-envelope block length \(q_a\) when applying
the tuning consequence below. Let \(\mathcal P\) be the
compatible nested cubic period family: every fixed power
of \(R\) divides all sufficiently late periods in
\(\mathcal P\).

For a \(q\)-step block with final period \(N\in\mathcal P\),
put
\begin{equation}
 M_j=MR^j,\qquad Q_j=R^{q-j}N,\qquad 0\le j\le q.
 \label{fbt-block-scales}
\end{equation}
Here \(M\) is any permitted preceding homogeneous depth.
Let \(D_{M_j,R;Q_{j+1}}\) be the actual finite leading
correction \(\Delta+\mathcal Q_{\det}\) of step \(j\).
Define
\begin{equation}
 b_{0,j}^{M,N}
    =\beta_{M_jR,Q_{j+1}}(D_{M_j,R;Q_{j+1}}),
 \qquad
 B_q(M,N)=\sum_{j=0}^{q-1}b_{0,j}^{M,N}.
 \label{fbt-finite-coefficients}
\end{equation}
Each functional is the actual finite short-lift ratio
normalized by its actual finite classical reference.
These coefficients are real on the real branches.
No finite ratio is replaced by an infinite moment.

The uniform outer-domain leading-family bound and the
finite moment estimate give, above the fixed reference
threshold,
\begin{equation}
                          |b_{0,j}^{M,N}|\le B_0
 \label{fbt-leading-bound}
\end{equation}
uniformly in depth and permitted volume. The reference
denominators are bounded below by \(1/2\).

\subsection{Varying-depth structural limits and their moments}
\label{fbt-compactness}

Take any \(M_n\) and any \(N_n\in\mathcal P\) tending to
infinity. The old depths need not be bounded or converge.
Extract jointly the classical references and background
maps at the shifted depths \(M_nR^j\), their required
fixed-order vertices, the actual retained flat precisions
and covariances, and all \(q\) leading corrections in
\eqref{fbt-finite-coefficients}.
The uniform strict support, common-domain and exponential
source-Hessian bounds provide such a subsequence. If a
structural reference comparison requires additional fixed
depth shifts or endpoint powers, include their countable
collection in the same diagonal extraction. Cofinal
divisibility guarantees their eventual finite availability.

Write the limiting corrections as \(D_j^*\), with true
infinite Maxwell functional \(\beta_{j+1}^*\) at output
scale \(j+1\), and set
\begin{equation}
                       B_q^*=\sum_{j=0}^{q-1}
                                         \beta_{j+1}^*(D_j^*).
 \label{fbt-limit-sum}
\end{equation}
The scale indices describe the structural limit. They
do not identify it with one fixed finite-depth system.

These are actual determinant structural limits. At fixed
\(R\) the retained fiber per current block has fixed finite dimension;
the whole retained operator still has dimension growing with volume.
Local matrix convergence and exponential envelopes pass
products and the inverse identity \(A_0C_0=I\) by truncating
the intermediate sums. The chart and pivot maps are fixed;
common-domain Cauchy estimates give convergence of classical
vertices through order four. The absolutely convergent
trace formula therefore identifies the limiting Hessian
with the same actual determinant expression. One may
equivalently extract the strict scalar families jointly
and identify their Hessians in this way. No old fine
coordinate count occurs.

\begin{lemma}[Passage of the actual finite moments]
\label{fbt-moment-lemma}
Along every such joint subsequence,
\begin{equation}
 b_{0,j}^{M_n,N_n}\longrightarrow
                      \beta_{j+1}^*(D_j^*)\quad(0\le j<q),
 \qquad
 B_q(M_n,N_n)\longrightarrow B_q^*.
 \label{fbt-moment-limits}
\end{equation}
\end{lemma}
\begin{proof}
For a fixed \(j\), the leading source-Hessian kernels
converge locally and have a common exponential row
bound. Restrict their short-lift second moment first
to a fixed displacement radius. The finite sum converges.
The rest has a uniformly summable exponential
second-moment tail. The finite short-lift cutoff tends
to infinity since \(Q_{j+1}=R^{q-j-1}N_n\to\infty\).

The actual finite reference denominator tends uniformly
to one by the flat normalized classical periodization
estimate, even for these varying depths. Division proves
the first assertion. The sum has fixed finite length
\(q\), proving the second. Neither an extensive total
limit nor a regulator convergence rate was used.
\end{proof}

\subsection{The conditional positive transfer}
\label{fbt-positive-section}

\paragraph{Unproved structural positivity hypothesis.}
Assume there is one \(\delta>0\) such that \emph{every}
joint structural infinite limit obtainable above satisfies
\begin{equation}
                                  B_q^*\ge\delta.
 \label{fbt-positive-hypothesis}
\end{equation}
The same margin must work for all varying-depth and
permitted-period sequences. Positivity in one selected
fixed-depth limit is insufficient.

\begin{proposition}[Uniform finite block lower bound]
\label{fbt-transfer-proposition}
Under \eqref{fbt-positive-hypothesis}, for every
\(0<\varepsilon<\delta\) there is a threshold
\(N_*(\varepsilon)\), independent of preceding depth,
such that
\begin{equation}
 B_q(M,N)\ge\delta-\varepsilon
 \quad\hbox{for all permitted }M,\quad
 N\in\mathcal P,\quad N\ge N_*(\varepsilon).
 \label{fbt-finite-lower}
\end{equation}
In particular the finite block is at least \(\delta/2\)
above one uniform period threshold.
\end{proposition}
\begin{proof}
Otherwise choose \(N_n\to\infty\) in \(\mathcal P\)
and old depths \(M_n\) such that
\(B_q(M_n,N_n)<\delta-\varepsilon\).
Take a joint structural subsequence. By
Lemma~\ref{fbt-moment-lemma}, its limiting sum is at most
\(\delta-\varepsilon\), contradicting
\eqref{fbt-positive-hypothesis}.
\end{proof}

The threshold is also uniform in a block's starting
stage: absorb that stage into its preceding depth \(M\).
Its final period is the \(N\) in
\eqref{fbt-block-scales}. The proof gives no numerical
threshold or convergence rate.

\subsection{The finite tuning input needs only this qualitative bound}
\label{fbt-tuning-section}

Take \(q\) to be a multiple of \(q_a\). Write
\(\mathcal F_{j,g}\) for the actual finite quantum extraction map,
\(x_j\) for the actual real extracted input class, and
\(d_j(g,x)=\beta_{j+1}(\mathcal F_{j,g}(x))\) for its
inverse-coupling decrement. On the established
finite activity envelope the actual decrement obeys
\begin{align}
 |d_j(g,x_j)-b_{0,j}|
 \le{}&C_{\rm int}\omega(N_{\min})
       +C_{{\rm par},\beta}
                    \{\tau(T(g))+g^2T(g)^6\}\notag\\
 &+C_{{\rm in},\beta}\,gT(g)/b.
 \label{fbt-actual-error}
\end{align}
The first term tends qualitatively to zero as the
period threshold increases. The other terms tend to
zero in the weak-coupling limit at the fixed field
radius and chosen real threshold \(T(g)\).

Use \eqref{fbt-finite-lower} with
\(\varepsilon=\delta/2\). Raise the period threshold
and lower the weak-coupling ceiling until
\eqref{fbt-actual-error} is at most \(\delta/(4q)\).
Then every permitted \(q\)-step block satisfies
\begin{equation}
 \frac{\delta}{4}\le
       \sum_{j=k}^{k+q-1}d_j(g_j,x_j)
       \le qD_{\rm abs}.
 \label{fbt-actual-block}
\end{equation}
The upper bound is the proved absolute-increment bound.
Repeating the \(q_a\)-step activity blocks gives the same
intermediate envelope throughout this longer block.
No new large-endpoint field-dilation estimate is used.

These are the positive block bounds used in the
continuous clipped-recursion shooting argument.
Put \(\Gamma_{\rm w}=g_{\rm w}^{-2}\), where
\(g_{\rm w}\) is the chosen weak-coupling ceiling.
For a terminal inverse coupling
\(\Gamma_*\ge\Gamma_{\rm w}+qD_{\rm abs}\), that argument
then gives an initial value in
\begin{equation}
 [\,\Gamma_*+m\delta/4,\ \Gamma_*+mqD_{\rm abs}\,]
 \label{fbt-shooting-interval}
\end{equation}
which makes \(mq\) genuine weak local extraction steps
and reaches \(\Gamma_*\). This remains conditional on
\eqref{fbt-positive-hypothesis}.

If an averaged formulation for arbitrary finite
interval lengths is desired, write \(n=mq+r\),
\(0\le r<q\), and use \eqref{fbt-leading-bound} on
the final fringe. For intervals entirely above the
common period threshold,
\begin{equation}
 \sum_{j=k}^{k+n-1}b_{0,j}
 \ge cn-C_{\rm fin},\qquad
 c=\frac{\delta}{2q},\qquad
 C_{\rm fin}=(q-1)(B_0+c).
 \label{fbt-optional-fringe}
\end{equation}
Indeed the complete blocks contribute at least
\(m\delta/2\), and the fringe contributes at least
\(-rB_0\). The shooting proof needs only the fixed-block
bound, so this additional formulation is optional.

\subsection{The remaining sign problem}
\label{fbt-scope}

Include the one long-ratio \(R^q\) determinant in the
joint extraction. The proved leading cocycle and true
infinite classical marginal preservation give
\begin{equation}
                         B_q^*
                          =\beta_q^*(D_{0,R^q}^*).
 \label{fbt-long-ratio}
\end{equation}
A uniformly positive reference comparison for this
coefficient in every structural limit would therefore
supply \eqref{fbt-positive-hypothesis}. The equality
itself proves no sign.

Once that positive structural margin is established,
qualitative compactness is enough for the needed finite
block lower bound; a separate exponential regulator
rate is unnecessary. This section does not establish
the margin, exact finite marginal additivity, summability
of errors over an infinite history, uniqueness of the
old-depth limit, or convergence of tuned local densities
to a continuum measure.
